\chapter{NDK Runtime Model}

The Nova Developer Kit is a ca65 assembly library stack over Nova's
memory-mapped hardware and NovaHost services. It is deliberately thin. Most
helpers name registers, copy arguments into the hardware mailbox a subsystem
expects, issue a command, and return a small status code.

\begin{center}
\novasvg[width=0.78\textwidth]{build/diagrams/ndk-runtime-stack}
\end{center}

\section{How Libraries Are Used}

Most libraries have an interface file and an implementation file. Include the
\novacode{.inc} file near the top of an application so symbols and
pseudo-register aliases are visible. Include or link the \novacode{.s} file
once after the code that calls it.

\begin{lstlisting}
.setcpu "65c02"

.include "fio.inc"
.include "vgc.inc"

main:
      JSR   fio_load
      JSR   vgc_vsync
      RTS

.include "fio.s"
.include "vgc.s"
\end{lstlisting}

Ordering matters for libraries that support size stripping. \novacode{fio.s}
and \novacode{audio.s} emit only referenced public routines by default, plus the
dependencies those routines need. If the implementation is included before the
call sites, ca65 has not seen the references yet and the routines can be
stripped. Define \novacode{NOVA\_EMIT\_ALL\_RUNTIME},
\novacode{FIO\_EMIT\_ALL\_RUNTIME}, or \novacode{AUDIO\_EMIT\_ALL\_RUNTIME}
for ROM and debug builds that intentionally keep full implementations.

\section{Calling Convention}

The runtime libraries share a small convention, but individual routines can
override it in their API entry:

\begin{itemize}
\item A return value of A=0 means success. A return value of A=1 usually means a
generic error, with detail in a library result byte or hardware error register.
\item A, X, and Y are scratch unless the routine documents them as outputs.
\item \novacode{NVR0L} through \novacode{NVR7H} are shared pseudo-register
argument bytes. They are not persistent storage.
\item Hardware-facing routines often leave detailed state in registers such as
\novacode{FIO\_STATUS}, \novacode{FIO\_ERRCODE}, \novacode{XMC\_STATUS}, or VGC
pending IRQ registers.
\end{itemize}

Treat pseudo-registers and hardware command registers as mailboxes. Load them
immediately before the call that needs them and copy out any result before the
next subsystem reuses the same bytes.

Routines whose names begin with \novacode{I\_} use Nova's inline-parameter ABI.
Their immutable argument bytes follow the \novacode{JSR}; the routine consumes
those bytes and advances the saved return address before returning. This keeps
static strings and their calls together without pointer setup. Inline blocks
must match the exact layout documented by the routine and are not ordinary
Pascal-callable NDK functions.

\begin{lstlisting}
      JSR   I_VTEXT_PRINT_AT
      .byte 12, 4, "Ready", 0
\end{lstlisting}

\section{Hardware And Host Boundaries}

\novacode{nova.inc} is the shared map for Nova hardware addresses and command
constants. VGC, DMA, blitter, XMC, NIC, timer, math, and FIO registers are all
memory mapped from the CPU's point of view, but they do not all behave like
RAM.

VGC memory spaces are off-bus display memories reached through VGC commands,
the VRAM port, FIO graphics transfers, the pager, DMA, or the blitter. XRAM is a
flat expansion memory device reached through XMC windows, XRAM helpers, DMA, or
FIO streaming. FIO is a host-command mailbox; NovaHost resolves mounted devices,
host files, media playback, runtime ROM loading, and random-number requests.

\section{Generated Reference}

The API chapter is generated from \novacode{@label} blocks in
\novacode{software/runtime/asm}. If an API box is wrong, the fix belongs in the source
comment next to the implementation, not in the generated LaTeX. Authored
chapters in this book should explain workflows and caveats; generated entries
should stay close to exact symbols, inputs, outputs, and required registers.
